% ==============================================================================
% T0 Theory / FFGFT: Document 275 (Chapter Content, EN)
% Title: From xi to phi -- Recursive Scale Hierarchy and Spectral Gap Structure
% Author: Johann Pascher
% Date: June 2026; revised 11 June 2026 (audit P. Stoychev, IPI)
% References: Doc. 009 (P10), 188, 190 (P35 triviality / P37 phi roles), 258, 270, 271, 272, 274; Krueger (HLV) Zenodo 20615649
% ==============================================================================

\newtheorem*{observationstar}{Observation}

\section*{Abstract}
We study the recursive damping sequence
$r(k{+}1)=r(k)(1-\xi)$, $r(0)=D_f/3$, which follows solely from the
fundamental geometric parameter $\xi = 4/30000$ and the fractal
dimension $D_f = 3-\xi$. At recursion depth
$k^* = \log(\varphi)/\xi \approx 3609$ the sequence \emph{passes
through} the value $1/\varphi$: $r(3609) \approx 1/\varphi$ with
absolute error $4{\cdot}10^{-5}$ (relative $6.5{\cdot}10^{-5}$); the
exact real $k^*_\mathrm{exact}=3608.51$ hits $1/\varphi$ identically.
\textbf{Status clarification (revision 11~June 2026):} $1/\varphi$ is
\emph{not a fixed point} of the recursion (the unique fixed point is
$0$) but a passage point; since $k^*(c)$ exists for \emph{every}
target $c\in(0,1)$ with identical parametric status, the
$\xi^N/\varphi^N$ triviality guardrail (Doc.~190 P35) applies to this
result in full. It would acquire physical content only through an
independent fixing of $k^*\approx 3609$ --- currently open.
The recursion levels are geometric scale levels with the FFGFT-internal
factor $q = 2/3$; $(2/3)^{3609}$ is convergent.
The numerical validation study was fully corrected (v4) after an
independent audit (P.~Stoychev, IPI, 11~June 2026): with a genuine 4D
implementation of the $T^4$ graph and dimension-matched 4D null models,
FFGFT~$T^4$ lies \emph{inside} the null band; HLV, with a genuine
ensemble, also lies inside the band. Earlier contrary findings (v3)
were implementation artefacts and are visibly retracted in this
document.

% ============================================================================
\section{Motivation and Context}
% ============================================================================

Doc.~258 showed that the Koide scalar
$Q_\mathrm{FFGFT} \approx 0.6677$ follows from the geometric invariants
$r_e, r_\mu, r_\tau$ and the lepton exponents $p_e, p_\mu, p_\tau$.
This value is an \emph{output} of the theory, not an input parameter.

In parallel, the HLV lattice construction of Krüger~\cite{krueger275en}
uses the golden ratio $\varphi = (1+\sqrt{5})/2$ as the projection ratio
in the 6D$\to$3D cut-and-project construction. Krüger's Run~E identifies
the low-sector gap hierarchy as the only remaining candidate separation
channel between HLV and null models.

This document answers three questions:
\begin{enumerate}
  \item Which scalar is physically correct for the macroscopic geometric
        structure?
  \item How does it connect to $\xi$?
  \item What does $k^* \approx 3609$ mean physically --- and why is it
        not an astronomical scale?
\end{enumerate}

\begin{remark}[All three FFGFT scalars are derivable]
The question is not which of the three scalars $\xi$, $1/\varphi$,
$Q_\mathrm{FFGFT}$ \emph{can} be derived from FFGFT --- that holds for
all three. The question is which one \emph{operates} at which physical
level and is therefore responsible for which observables.
\end{remark}

% ============================================================================
\section{Fundamental Parameters and Scalar Separation}
% ============================================================================

FFGFT works with the following parameter-free fundamental quantities:
\begin{align}
  \xi   &= \frac{4}{30000} \approx 1.333 \times 10^{-4}
    \qquad \text{(fractal deviation, microscopic)} \label{eq:xi275en}\\
  D_f   &= 3 - \xi \approx 2.99987
    \qquad \text{(Hausdorff dimension of }T^4\text{)} \label{eq:Df275en}\\
  q     &= \frac{2}{3}
    \qquad \text{(geometric recursion factor)} \label{eq:q275en}
\end{align}

\begin{definition}[Three FFGFT scalars near 0.6]\label{def:scalars275en}
\begin{center}
\begin{tabular}{@{}llp{0.44\linewidth}@{}}
\toprule
Quantity & Value & Role \\
\midrule
$\xi = 4/30000$ & $1.333\times10^{-4}$ &
  fundamental input; fractal dimension; active everywhere \\
$1/\varphi$ & $0.6180$ &
  macroscopic shell ratio; passage point of the $\xi$-recursion (output; status see Sect.~\ref{sec:passage275en}) \\
$Q_\mathrm{FFGFT}$ & $0.6677$ &
  Koide scalar leptons (output, Doc.~258); particle physics scale only \\
\bottomrule
\end{tabular}
\end{center}
All three are derivable from FFGFT. They are \emph{not} interchangeable.
\end{definition}

% ============================================================================
\section{The $\xi$-Damping Recursion}
% ============================================================================

\subsection{Definition and Closed Form}

\begin{definition}[$\xi$-damping recursion]\label{def:rec275en}
Let $r(0) = D_f/3 = (3-\xi)/3$ be the shell ratio at the lowest scale level.
The recursion reads:
\begin{equation}
  r(k+1) = r(k)\cdot(1-\xi), \qquad k = 0,1,2,\ldots
  \label{eq:rec275en}
\end{equation}
Closed form:
\begin{equation}
  r(k) = \frac{D_f}{3}\cdot(1-\xi)^k \approx \frac{D_f}{3}\cdot e^{-\xi k}.
  \label{eq:closed275en}
\end{equation}
\end{definition}

The physical content: $r(0) = D_f/3$ is the ratio of the fractal
dimension to the topological dimension~3. Each scale level damps this
ratio by the factor $(1-\xi)$.

\subsection{Main Observation (Revised): $1/\varphi$ as Passage Point}\label{sec:passage275en}

\begin{remark}[Revision 11~June 2026]
The original version of this section called $1/\varphi$ a \emph{fixed
point} of the recursion and stated the result as a theorem. Both were
wrong or overstated: the unique fixed point of $r\mapsto r(1-\xi)$ is
$r=0$; $1/\varphi$ is a \emph{passage point}. The correction follows
the audit by P.~Stoychev (IPI list, 11~June 2026).
\end{remark}

\begin{observationstar}
The sequence $r(k)$ passes through the value $1/\varphi$ at recursion
depth
\begin{equation}
  k^* = \frac{\log\varphi}{\xi}
      = \frac{\log\!\bigl(\tfrac{1+\sqrt{5}}{2}\bigr)}{4/30000}
      \approx 3609,
  \label{eq:kstar275en}
\end{equation}
with $r(3609) = 0.6179940013\ldots$ and
$|r(3609)-1/\varphi| = 4.0\times10^{-5}$
(relative $6.5\times10^{-5}$). The exact real
$k^*_\mathrm{exact} = \bigl(\log(1/\varphi)-\log(D_f/3)\bigr)/\log(1-\xi)
= 3608.51$ satisfies $r(k^*_\mathrm{exact}) = 1/\varphi$ identically.
\end{observationstar}

\begin{remark}[Parametric status --- the decisive caveat]
The arithmetic is correct, but the statement is \emph{not}
distinguished: for \emph{every} target $c\in(0,1)$ there exists
\begin{equation}
  k^*(c) = \frac{\log c - \log(D_f/3)}{\log(1-\xi)}
\end{equation}
with exactly the same parametric status --- \enquote{no free
parameter} holds for all targets alike, because the free exponent has
merely been renamed $k^*$ and solved for the target:
\begin{center}
\begin{tabular}{@{}lr@{}}
\toprule
Target $c$ & $k^*(c)$ \\
\midrule
$1/\varphi = 0.6180$        & $3608.5$ \\
$Q_\mathrm{FFGFT} = 0.6677$ & $3028.8$ \\
$1/2$                        & $5197.9$ \\
$1/e$                        & $7499.2$ \\
$\alpha_\mathrm{em}$        & $36899.0$ \\
\bottomrule
\end{tabular}
\end{center}
The $\xi^N/\varphi^N$ triviality guardrail (Doc.~190 P35) therefore
applies to this very result. The observation would acquire physical
content only if $k^*\approx 3609$ were fixed by something independent
(a level count of the construction, an observable) and \emph{then}
landed on $1/\varphi$. No such fixing is currently known.
\end{remark}

% ============================================================================
\section{Scale Interpretation: $q=2/3$, Not Factor~2}
% ============================================================================

\begin{remark}[Critical point]\label{rem:scale275en}
A naive reading might interpret $k^* \approx 3609$ as astronomically many
renormalisation-group steps: with factor~2 per step that would be a scale
spread of $2^{3609} \approx 10^{1086}$, far beyond any experimental reach.

This objection misses the FFGFT structure.
\end{remark}

The four torus lengths follow $L_\mu = L_0\cdot q^\mu$ with $q = 2/3$.
The cumulative scaling after $k$ levels is $S(k) = (2/3)^k \to 0$.

\begin{proposition}[Convergence with FFGFT scaling]\label{prop:conv275en}
Comparison:
\begin{center}
\begin{tabular}{@{}lrr@{}}
\toprule
Interpretation & Scale factor at $k^*=3609$ & Character \\
\midrule
RG, factor~2 (external convention) & $2^{3609}\approx 10^{1086}$ & divergent \\
FFGFT, $q=2/3$ (internal)          & $(2/3)^{3609}\approx 10^{-635}$ & convergent \\
\bottomrule
\end{tabular}
\end{center}
$k^* \approx 3609$ is a natural number within the theory's scale spectrum,
not an astronomical limit.
\end{proposition}

% ============================================================================
\section{Spectral Validation Study (Three Layers)}
% ============================================================================

\subsection{Design}

\textbf{Script version v4 (revised after audit).} Two separate
sectors are compared, each against dimension-matched null models:
\emph{4D sector} --- FFGFT~$T^4$ as a kNN graph in genuine 4D
coordinates (anisotropic torus lengths $L_\mu = 2\pi(2/3)^\mu$,
anisotropy $L_0/L_3 = 3.38$) against three 4D nulls (isotropic cubic,
jittered anisotropic, random anisotropic box); \emph{3D sector} ---
HLV G-Lattice ($k$NN graph, now with a genuine seed ensemble
$\sigma=0.05$) against three 3D nulls (jittered cubic, random, cubic).
$N=2000$, $k=16$, three seeds for all models; uniqueness of all point
sets is enforced by assertion.

\begin{remark}[What was wrong in v3 --- audit P.~Stoychev]
(i)~The v3 \enquote{$T^4$ graph} used the grid coordinate $n_3$ only
in the acceptance norm, not in the position: of 2000 points only 729
were unique 3D positions; the $1/d^2$ weighting with clamp produced
Laplacian eigenvalues up to $2.5\times10^{12}$.
(ii)~\texttt{hlv\_points(seed)} never used its seed --- the
\enquote{ensemble} was the same computation three times.
(iii)~The shell detector was structurally biased for targets $<1$
(see Sect.~\ref{sec:layer2-275en}). All three errors are fixed in v4;
in addition, 4D models are compared only against 4D nulls, since
comparing a 4D graph against 3D nulls would again be unfair.
\end{remark}

\begin{description}
  \item[Layer~1] Gap-hierarchy metrics: max.\ gap, gap-CV,
    low-sector $\langle r\rangle$, IDS --- analogous to Krüger Run~E.
  \item[Layer~2] Shell structure: $1/\varphi$ vs.\ $Q_\mathrm{FFGFT}$
    on the HLV point cloud at $N=2000$.
  \item[Layer~3] Numerical verification of $k^*$ and $r(k^*)$.
\end{description}

\begin{remark}[Heuristics in the script]
The shell-structure function uses a 5\% threshold for shell detection.
This is a heuristic; for a publication the radial distribution function
$g(r)$ would be more robust. The ensemble size (3~seeds) is illustrative,
not statistically secured.
\end{remark}

\subsection{Results Layer~1 (v4): Gap Hierarchy ($N=2000$, 3~Seeds)}

\begin{remark}[Retraction of the v3 result]
The result reported in the first version
(\enquote{FFGFT $T^4$ graph outside the null band,
gap\_CV $=1.68\pm0.09$}) is \textbf{retracted}: it was entirely an
artefact of point duplicates (729 unique of 2000 points) combined with
the $1/d^2$ clamp. The precision statement
\enquote{HLV $\approx$ jittered, $\Delta<0.001$ across 3 seeds} is
also withdrawn --- it compared a single realisation with itself
(seed bug).
\end{remark}

\begin{table}[ht]
\centering
\caption{Gap CV (low sector, $n{=}100$ modes, $N{=}2000$,
mean~$\pm$~std over 3~seeds), v4: separate sectors with
dimension-matched null models.}
\label{tab:gap275env4}
\begin{tabular}{@{}llrl@{}}
\toprule
Sector & Model & gap\_CV & Band status \\
\midrule
4D & \textbf{FFGFT $T^4$ (anisotropic)} & $1.329\pm0.025$ & in band \\
4D & Null: isotropic cubic              & $2.406\pm0.010$ & band edge (top) \\
4D & Null: jittered anisotropic         & $0.690\pm0.026$ & band edge (bottom) \\
4D & Null: random anisotropic box       & $0.740\pm0.037$ & -- \\
\midrule
3D & \textbf{HLV G-Lattice} ($\sigma{=}0.05$) & $0.710\pm0.008$ & in band \\
3D & Null: jittered cubic               & $0.768\pm0.011$ & -- \\
3D & Null: random                       & $0.675\pm0.043$ & band edge (bottom) \\
3D & Null: cubic                        & $0.978\pm0.010$ & band edge (top) \\
\bottomrule
\end{tabular}
\end{table}

Findings:

\begin{enumerate}
  \item \textbf{No model separates.} FFGFT~$T^4$ lies at $1.329$
        inside the 4D null band $[0.690,\,2.406]$; HLV lies at $0.710$
        inside the 3D null band $[0.675,\,0.978]$.
  \item \textbf{The isotropic 4D lattice itself has the highest
        gap\_CV} ($2.406$) --- regular lattices produce strong
        eigenvalue degeneracies in 4D. A high gap\_CV is therefore not
        a $T^4$ signature but a generic order signature.
  \item \textbf{Robustness checked:} the $T^4$ result is
        jitter-stable ($\sigma=0.02$--$0.10$: $1.43/1.35/1.35$) and
        consistent for $n_\mathrm{range}=4/5$ ($1.35/1.40$);
        $n_\mathrm{range}=3$ ($2.28$) is a small-size artefact.
  \item \textbf{HLV with a genuine ensemble} remains in the band and
        close to jittered cubic --- the v3 \emph{finding} (consistent
        with Krüger's Run~E) survives, the precision claimed then does
        not.
\end{enumerate}

\begin{remark}[Consequence for the comparison]
At $N=2000$ the gap\_CV diagnostic separates neither the $T^4$ graph
nor HLV from dimension-matched null models. This \emph{reinforces}
Krüger's conservative Run-E finding from an independent pipeline and
further constrains the common-spectral-signature hypothesis of
Doc.~271: at this size, no signature is demonstrable at the graph
level.
\end{remark}

\subsection{Results Layer~2: Shell Structure --- Retracted}\label{sec:layer2-275en}

\begin{remark}[Retraction of the v3 discrimination]
The discrimination reported in the first version
($\mathrm{err}(Q_\mathrm{FFGFT}) < \mathrm{err}(1/\varphi)$,
$\Delta = 0.145$) is \textbf{retracted}. The detector identifies
shells as a strictly increasing distance sequence (factor
$\geq 1.05$); for any target ratio $c<1$ the minimum achievable error
is $(1.05/c - 1)$: $0.699$ for $1/\varphi$, $0.573$ for
$Q_\mathrm{FFGFT}$. \emph{The larger of two sub-1 ratios wins on any
point set whatsoever} --- the script's own controls showed this
(cubic: $0.976$ vs.\ $1.135$; random: $1.194$ vs.\ $1.371$) but were
not read as an alarm. The finding was structural, not physical
(audit P.~Stoychev).
\end{remark}

In v4 the detector raises an exception for targets $\leq 1$; only
ratios $>1$ are tested. Reference value:
$\mathrm{err}(\varphi{=}1.618)_\mathrm{HLV} = 0.257$; the best-fit
consecutive shell ratio of the HLV point cloud is, per Stoychev's
measurement, $\approx 1.15$ --- neither $\varphi$, $1/\varphi$, nor
$Q_\mathrm{FFGFT}$. For future versions, the radial distribution
function $g(r)$ with peak-position analysis is the more robust route.

\subsection{Result Layer~3: Numerical Verification}

\begin{equation}
  k^* = \frac{\log\varphi}{\xi} \approx 3609.1,
  \qquad
  r(3609) = 0.6179940013\ldots \approx \frac{1}{\varphi}
  \quad (|\text{error}| = 4.0\times10^{-5},
  \text{ relative } 6.5\times10^{-5}).
\end{equation}

The arithmetic is correct; the error statement was made precise in v3
(previously incorrectly $<10^{-9}$). For the \emph{status} of the
statement --- passage point, not fixed point; $k^*(c)$ exists for every
target --- see Sect.~\ref{sec:passage275en}: \enquote{without free
parameters} does not distinguish this result.

% ============================================================================
\section{Discussion}
% ============================================================================

\subsection{Relation to the $\varphi^N$-triviality warning (Doc.~190 P35) --- Revised}

\begin{remark}[Retraction of the demarcation]
The first version of this section gave three arguments why the
triviality warning would \emph{not} apply to the $k^*$ result
(\enquote{no free exponent}, \enquote{no empirical anchor},
\enquote{no external calibration}). This demarcation is
\textbf{retracted}: all three arguments overlook that the free
exponent has merely been \emph{renamed} $k^*$ and \emph{solved for}
the desired target. $k^*(c)$ exists for every $c\in(0,1)$ with
identical properties (table in Sect.~\ref{sec:passage275en}); the
selection of the target $1/\varphi$ is the hidden choice. P35 applies
in full --- this is exactly the case the guardrail was written for,
and applying it to one's own result is its first serious
corpus-internal stress test (passed: the result is downgraded).
\end{remark}

What remains: the three-level separation of $\varphi$-roles documented
in Doc.~190 P37 --- (1)~microscopic/$\xi$-level,
(2)~macroscopic/$1/\varphi$ as passage point of the recursion,
(3)~particle physics/$Q_\mathrm{FFGFT}$ --- remains valid and useful
as a \emph{conceptual separation}. Doc.~188 uses $\varphi^{-1}$ as a
phenomenological scaling factor between fractal levels; whether this
factor follows structurally from the $\xi$-recursion is, after the
present revision, again an \emph{open question}: the recursion
delivers it only if the depth $k^*\approx3609$ is independently
justified.

\subsection{Position in the Document Series}

Doc.~270 describes the reduction chain $T^4\to T^0$ as two qualitatively
different operation types. The present document adds the \emph{scale
hierarchy}: the $q=2/3$ recursion connects the microscopic $\xi$-level
with the macroscopic $\varphi$-level through cumulative geometric
damping --- without a qualitative change in operation type.

Doc.~271 provides a structural comparison of FFGFT and HLV. The present
document makes the spectral comparison announced there concrete through
the numerical script and identifies gap\_CV as a shared diagnostic channel.

\subsection{Open Points}

\begin{enumerate}
  \item Shell diagnostics on a $g(r)$ basis (peak positions, ratios
        $>1$) --- the v3 detector is unusable for sub-1 targets
        (Sect.~\ref{sec:layer2-275en}).
  \item Pre-declared ensemble band tests analogous to Krüger Run~E
        (metrics, bands and seeds fixed before the computation) for
        the 4D $T^4$ graph.
  \item Larger $N$ and more seeds: the present 3-seed ensembles are
        illustrative, not statistically secured.
  \item An independent fixing of $k^* \approx 3609$ (a level count of
        the construction or a physical observable) --- without it the
        passage-point observation has no physical content
        (Sect.~\ref{sec:passage275en}). Currently not known.
  \item Cross-validation of the v4 implementation against the
        deduplicated/seeded variants offered by P.~Stoychev.
\end{enumerate}

% ============================================================================
\section{Conclusion}
% ============================================================================

\begin{enumerate}
  \item The $\xi$-damping recursion passes through $1/\varphi$ at
        $k^* = \log(\varphi)/\xi \approx 3609$; the arithmetic is
        correct (error $4{\cdot}10^{-5}$), but $1/\varphi$ is a
        \emph{passage point, not a fixed point}, and $k^*(c)$ exists
        for every target --- the P35 triviality guardrail applies;
        without an independent fixing of $k^*$ the observation has no
        physical content.
  \item With $q = 2/3$, $(2/3)^{3609}$ is convergent and internally
        consistent, not an astronomical limit --- this scale
        interpretation stands.
  \item \textbf{v4 main result ($N=2000$, dimension-matched nulls):}
        No separation. FFGFT~$T^4$ (genuine 4D,
        gap\_CV $=1.329\pm0.025$) lies inside the 4D null band
        $[0.690,\,2.406]$; HLV (genuine ensemble, $0.710\pm0.008$)
        lies inside the 3D null band $[0.675,\,0.978]$. This
        reinforces Krüger's conservative Run-E finding from an
        independent pipeline.
  \item The v3 findings \enquote{$T^4$ outside the band} and
        \enquote{$Q_\mathrm{FFGFT}$ discriminates against
        $1/\varphi$} are \textbf{retracted} (duplicate artefact and
        detector bias, respectively); the HLV ensemble seed bug is
        fixed. All three errors trace back to the audit by
        P.~Stoychev and are registered as K4 in Doc.~190.
  \item The three scalars $\xi$, $1/\varphi$, $Q_\mathrm{FFGFT}$
        remain useful as a conceptual separation at different physical
        levels (Doc.~190 P37); the \emph{derivation claim} for
        $1/\varphi$ from the recursion is downgraded to an open
        question (Doc.~190 P38).
\end{enumerate}

\paragraph{Connections.}
Doc.~258 ($Q_\mathrm{FFGFT}$, Koide ladder), Doc.~270 (reduction chain
$T^4\to T^0$), Doc.~271 (FFGFT-HLV structural comparison), Doc.~272
(HLV/Dot-Theory context). The Python script
\texttt{ffgft\_hlv\_gap\_v4.py} (corrected version; v2/v3 remain in
the archive as documented error states) is publicly available and
reproducible.

\paragraph{Acknowledgement.}
Plamen Stoychev (IPI) for the independent audit of 11~June 2026, which
identified all three script errors and forced the status clarification
of the $k^*$ result. The revision of this document implements his audit
in full.

\begin{thebibliography}{9}

\bibitem{krueger275en}
M.~Krüger,
\emph{Operational Falsification Pathways for the HLV G-Lattice},
Zenodo (2026),
\url{https://doi.org/10.5281/zenodo.20615649}.

\bibitem{doc258en}
J.~Pascher, \emph{Doc.\,258: Geometric ladder $q=2/3$},
FFGFT corpus (2026).

\bibitem{doc270en}
J.~Pascher, \emph{Doc.\,270: Dimensional change in FFGFT},
FFGFT corpus (2026).

\bibitem{doc271en}
J.~Pascher, \emph{Doc.\,271: FFGFT compared with HLV},
FFGFT corpus (2026).

\bibitem{doc190en}
J.~Pascher, \emph{Doc.\,190: Correction register FFGFT/T0
(K1--K4, P1--P38)}, FFGFT corpus (2026).

\bibitem{doc009en}
J.~Pascher, \emph{Doc.\,009: Origin of $\xi = 4/30000$},
FFGFT corpus (2026).

\bibitem{doc188en}
J.~Pascher, \emph{Doc.\,188: Geometric foundations of FFGFT},
FFGFT corpus (2026).

\end{thebibliography}
